\subsection{The Patch-Induced Frequency}

The tokeniser introduces a characteristic \emph{patch frequency}:
\begin{equation}
	f_{\mathrm{patch}} = \frac{f_s}{S}.
\end{equation}
This is the rate at which new patches are produced (in Hz, referred to original time),
and it plays a role analogous to a resampling rate in the token domain.
Intuitively, $f_{\mathrm{patch}}$ sets the \emph{temporal resolution} of the token
sequence: a smaller stride $S$ yields a higher patch frequency and thus a finer-grained
view of the signal's evolution, while a larger stride coarsens the view and reduces
computational cost.

\begin{example}[Computing $f_{\mathrm{patch}}$ for Chronos-Bolt]
	Chronos-Bolt uses $f_s = 512\,\text{Hz}$ and stride $S = 16$.
	Therefore:
	\[
	f_{\mathrm{patch}} = \frac{512}{16} = 32\,\text{Hz}.
	\]
	This means the model receives one new token every $1/32 \approx 31.25\,\text{ms}$,
	regardless of how fast the underlying signal actually varies.
	Any signal whose periodicity is commensurate with this $32\,\text{Hz}$ grid is at risk
	of phase-locking, as formalised below.
\end{example}

% -----------------------------------------------------------------
\subsection{Phase-Locking Condition}

Consider a sinusoidal signal $x[n] = A\cos\!\bigl(2\pi f n / f_s + \phi\bigr)$.

\begin{definition}[Phase-Locking]
	The signal \emph{phase-locks} to the patch grid when its period $T_0 = f_s/f$
	(in samples) is \emph{commensurate} with $S$:
	\begin{equation}
		\frac{S \cdot f}{f_s} \in \mathbb{Q}.
	\end{equation}
	The \emph{degenerate} (critical) case occurs when this ratio is an integer:
	\begin{equation}
		f = m \cdot \frac{f_s}{S} = m \cdot f_{\mathrm{patch}},
		\qquad m \in \mathbb{Z}^+,
	\end{equation}
	i.e.\ $f$ is a \textbf{harmonic of} $f_{\mathrm{patch}}$.
\end{definition}

The three regimes — integer ratio, rational ratio, and irrational ratio — produce
qualitatively different behaviours, summarised below.

\begin{example}[Three regimes of phase-locking]
	Let $f_s = 100\,\text{Hz}$ and $S = 10$, so $f_{\mathrm{patch}} = 10\,\text{Hz}$.
	
	\smallskip
	\noindent\textbf{Case 1 — Integer ratio (full degeneracy).}
	$f = 20\,\text{Hz}$, so $T_0 = 5$ samples and $S/T_0 = 2 \in \mathbb{Z}$.
	Every patch starts at exactly the same phase of the sinusoid; all patches are
	identical and the token sequence is constant.
	
	\smallskip
	\noindent\textbf{Case 2 — Rational ratio (partial degeneracy).}
	$f = 15\,\text{Hz}$, so $T_0 = 100/15 \approx 6.67$ samples and
	$S/T_0 = 3/2 \in \mathbb{Q} \setminus \mathbb{Z}$.
	Patches alternate between two distinct phase offsets: the sequence is periodic
	with period 2, not constant, so some information survives.
	
	\smallskip
	\noindent\textbf{Case 3 — Irrational ratio (no phase-locking).}
	$f = 10\sqrt{2} \approx 14.14\,\text{Hz}$, so $S \cdot f / f_s = \sqrt{2} \notin
	\mathbb{Q}$.
	The starting phase advances by an irrational fraction of $2\pi$ at each step;
	the patches never repeat and the full phase information is preserved.
\end{example}

% -----------------------------------------------------------------
\subsection{Effect of Patch Size $P$}

A single patch of $P$ samples spans a duration $P/f_s$ seconds and covers
$P \cdot f / f_s$ cycles of the signal.
When
\begin{equation}
	\frac{P \cdot f}{f_s} \in \mathbb{Z}^+,
	\quad \text{i.e.} \quad
	f = k \cdot \frac{f_s}{P}, \quad k \in \mathbb{Z}^+,
\end{equation}
the patch contains an exact integer number of cycles: its content is
\emph{internally redundant}.

\begin{example}[Within-patch redundancy]
	Let $P = 8$, $f_s = 8\,\text{Hz}$, $f = 2\,\text{Hz}$ ($T_0 = 4$ samples).
	A single patch contains $P \cdot f / f_s = 8 \cdot 2 / 8 = 2$ complete cycles:
	\[
	\mathbf{p} = [0, 1, 0, -1,\; 0, 1, 0, -1].
	\]
	The second half is an exact copy of the first half.
	The patch is internally redundant: it carries no more information than a patch of
	length $P/2 = 4$ would.
	Any linear feature extractor applied to this patch can at best recover the waveform
	shape, but cannot distinguish whether the signal has been observed for 4 or 8 samples.
\end{example}

The most severe degeneracy occurs when \emph{both} the stride condition and the patch-size
condition hold simultaneously:
\begin{equation}
	f = \frac{\ell \cdot f_s}{\mathrm{lcm}(P,\, S)},
	\qquad \ell \in \mathbb{Z}^+.
\end{equation}
At these frequencies, each patch is internally redundant \emph{and} all patches are
identical — the token sequence is maximally uninformative.

\begin{example}[Maximal degeneracy]
	Let $f_s = 100\,\text{Hz}$, $P = 6$, $S = 4$,
	so $\mathrm{lcm}(6,4) = 12$.
	The joint condition requires $f = \ell \cdot 100/12 \approx 8.33\,\ell\,\text{Hz}$.
	For $\ell = 1$: $f \approx 8.33\,\text{Hz}$, $T_0 = 12$ samples.
	Every patch spans exactly $6/12 = 0.5$ cycles (internally redundant in the sense
	that $T_0$ divides $\mathrm{lcm}(P,S)$), and the stride advances by exactly
	$4/12 = 1/3$ of a period, which is rational — triggering the partial phase-lock of
	Case 2 above and producing a two-periodic token sequence with minimal diversity.
\end{example}

% -----------------------------------------------------------------
\subsection{Token-Space Degeneracies}

\paragraph{At $f = m\,f_s/S$ (harmonics of $f_{\mathrm{patch}}$).}
\begin{itemize}
	\item All patches are translates of the same waveform by a phase that is a multiple
	of $2\pi$: they become \textbf{identical} (cf.\ Point~1).
	\item The token sequence is constant; it encodes \textbf{no phase information}.
	\item Two signals at the same frequency but different initial phases (shifted by any
	multiple of $S$ samples) map to the same token sequence:
	\emph{token-space invariance}.
\end{itemize}

\paragraph{At $f = k\,f_s/P$ (but not harmonics of $f_{\mathrm{patch}}$).}
\begin{itemize}
	\item Each patch is internally periodic, so its content is \emph{within-patch
		redundant}.
	\item Inter-patch variation still exists (partial degeneracy): the model retains
	some discriminative power, but each token is a compressed, information-poor
	representation of the signal at that time step.
\end{itemize}

% -----------------------------------------------------------------
\subsection{Overlap, Stride, and Blind-Spot Frequencies}

When $S < P$ (overlapping patches), consecutive patches share $P - S$ samples, partially
mitigating the aliasing effect even near harmonic frequencies.
Intuitively, the shared samples act as a \emph{bridge} between consecutive tokens,
preserving some phase continuity even when the stride would otherwise trigger
phase-locking.

\begin{example}[How overlap rescues phase information]
	Let $f_s = 100\,\text{Hz}$, $f = 25\,\text{Hz}$ ($T_0 = 4$ samples), $P = 6$,
	$S = 4$ ($= T_0$, so the stride condition is met).
	Without overlap ($S = P = 4$): $\mathbf{p}_0 = [0,1,0,-1]$ and
	$\mathbf{p}_1 = [0,1,0,-1]$ — identical.
	With overlap ($S = 4 < P = 6$):
	\[
	\mathbf{p}_0 = [x_0,x_1,x_2,x_3,x_4,x_5]
	= [0,1,0,-1,0,1],
	\]
	\[
	\mathbf{p}_1 = [x_4,x_5,x_6,x_7,x_8,x_9]
	= [0,1,0,-1,0,1].
	\]
	In this particular case the patches are still identical because $T_0 = 4$ also
	divides $P = 6$ (partially) — but for a generic $P$ not divisible by $T_0$, the
	shared samples would introduce a phase offset and break the degeneracy.
	This illustrates that overlap is a necessary, but not always sufficient, safeguard.
\end{example}

For $S \geq P$ (no overlap or gaps), the aliasing is maximal; the
\textbf{blind-spot frequencies} are:
\begin{equation}
	\boxed{f_{\mathrm{blind}} = m \cdot \frac{f_s}{S},
		\qquad m = 1, 2, 3, \ldots, \left\lfloor \frac{S}{2} \right\rfloor.}
\end{equation}
The upper limit follows from the Nyquist criterion applied to the token sequence, whose
effective sample rate is $f_s/S$.
These consequences follow directly from what is expressed in
\cite{paganiPreliminaryInsightsChronos2026b}.

\begin{example}[Blind-spot table for Chronos-Bolt]
	With $f_s = 512\,\text{Hz}$ and $S = 16$, we have $f_{\mathrm{patch}} = 32\,\text{Hz}$
	and $\lfloor S/2 \rfloor = 8$ blind spots:
	\begin{center}
		\begin{tabular}{@{}ccc@{}}
			\toprule
			$m$ & $f_{\mathrm{blind}}$ (Hz) & Interpretation \\
			\midrule
			1 & 32  & Fundamental patch frequency \\
			2 & 64  & 2nd harmonic \\
			3 & 96  & 3rd harmonic \\
			4 & 128 & 4th harmonic \\
			5 & 160 & 5th harmonic \\
			6 & 192 & 6th harmonic \\
			7 & 224 & 7th harmonic \\
			8 & 256 & Nyquist limit of token sequence \\
			\bottomrule
		\end{tabular}
	\end{center}
	Any input signal whose dominant frequency falls on one of these values will produce a
	degenerate token sequence and trigger a measurable drop in forecasting accuracy.
\end{example}

% -----------------------------------------------------------------
\subsection{Summary}

\begin{table}[h]
	\centering
	\begin{tabular}{@{}p{5.5cm}p{7.5cm}@{}}
		\toprule
		\textbf{Condition} & \textbf{Effect} \\
		\midrule
		$f = m\,f_s/S$
		& Full phase-locking $\to$ identical patches $\to$ constant token sequence. \\[4pt]
		$f = k\,f_s/P$
		& Internal patch periodicity $\to$ within-patch redundancy. \\[4pt]
		Both simultaneously
		& Maximal degeneracy: constant \emph{and} redundant token sequence. \\[4pt]
		$S < P$ (overlap)
		& Partial mitigation; residual phase information between tokens. \\[4pt]
		$S \geq P$ (no overlap)
		& Full degeneracy at blind-spot frequencies; no mitigation. \\
		\bottomrule
	\end{tabular}
	\caption{Summary of conditions and their effects on token-space observability.}
\end{table}
